\subsection*{Rank-one Jacobian characters force an actual CM quotient}

\paragraph{Statement.}
Let $J$ be an abelian variety over a number field $L$, let $\ell$ be a prime,
and let $W\subset V_\ell(J)$ be a nonzero one-dimensional
$\mathbf Q_\ell$-subspace stable under $G_L$.
There is an $L$-simple abelian variety $B$ and an actual surjective
$L$-homomorphism $q:J\to B$ such that $V_\ell(q)|_W\ne0$ and
\[
 \operatorname{End}^0_L(B)=E,\qquad [E:\mathbf Q]=2\dim B.
\]
Thus $B$ is of CM type, with its CM endomorphisms defined over $L$.
There is a place $\lambda\mid\ell$ of $E$ with $E_\lambda=\mathbf Q_\ell$;
the image of $W$ is the corresponding one-dimensional summand of
$V_\ell(B)$.

\paragraph{Proof with the projector rank retained.}
Choose an $L$-isogeny decomposition of $J$ into powers of $L$-simple
abelian varieties. Its rational inverse followed by one factor projection,
and then multiplication by a positive integer clearing denominators,
is an actual homomorphism $q:J\to B$ which is nonzero on $W$ for at least
one factor $B$. This map is surjective, since its image is a nonzero
abelian subvariety of the $L$-simple $B$. Denote its one-dimensional
image on $W$ by $W_B$.

Put $D=\operatorname{End}^0_L(B)$, a division algebra with centre $E$,
and write $[D:E]=d^2$. For a fixed embedding $L\hookrightarrow\mathbf C$,
the vector space $H_1(B(\mathbf C),\mathbf Q)$ is a left $D$-vector space
of an integer dimension $m\ge1$. Consequently
\[
 2\dim B=m d^2[E:\mathbf Q].
\]
The Betti--etale comparison preserves every endomorphism of $B$, so
$V_\ell(B)$ is, as a left $D\otimes\mathbf Q_\ell$-module, the direct
sum of $m$ copies of its regular module.

Faltings' semisimplicity and endomorphism theorem supply a $G_L$-equivariant
projector $e$ with image $W_B$ and identify it with an idempotent in
$D\otimes\mathbf Q_\ell$.
The exact original locators are \cite[\S5, S\"atze3--4,p.360]{GenFaltings1983};
the corresponding Hom theorem is Korollar1,p.361. In particular the
projector is allowed to have $\ell$-adic coefficients: it is not silently
asserted to descend to a rational idempotent of $D$.

For each $\lambda\mid\ell$, write
\[
 D\otimes_E E_\lambda\simeq M_{r_\lambda}(D_\lambda),\qquad
 [D_\lambda:E_\lambda]=\delta_\lambda^2,\qquad
 r_\lambda\delta_\lambda=d.
\]
An idempotent in this factor of rank $s_\lambda$ over the division
algebra $D_\lambda$ is conjugate to a diagonal matrix with
$s_\lambda$ entries equal to one. Its image on the left regular module
has $\mathbf Q_\ell$-dimension
\[
 s_\lambda r_\lambda\delta_\lambda^2[E_\lambda:\mathbf Q_\ell]
 =s_\lambda d\delta_\lambda[E_\lambda:\mathbf Q_\ell].
\]
It follows that the actual rank of $e$ on $V_\ell(B)$ is
\[
 1=m\sum_{\lambda\mid\ell}
 s_\lambda d\delta_\lambda[E_\lambda:\mathbf Q_\ell].
\]
Every summand is a nonnegative integer. There is exactly one nonzero
summand, and all its factors are one. In particular $m=d=1$,
$D=E$, $[E:\mathbf Q]=2\dim B$, and the active completion
$E_\lambda$ equals $\mathbf Q_\ell$. This proves the stated CM quotient
and the exact summand identification. No assertion that $\dim B=1$
has been used or obtained.

For completeness, $E$ is a CM field, not only a commutative algebra
of the required degree. Choose a polarization over $L$. Its Rosati
involution $\dagger$ preserves $E=\operatorname{End}^0_L(B)$.
The associated positive real inner product on $H_1(B(\mathbf C),\mathbf R)$
makes $a^\dagger$ the adjoint of $a$, so
\[
 \operatorname{Tr}_{E/\mathbf Q}(aa^\dagger)>0\qquad(a\ne0).
\]
Here the equality with the ordinary representation trace uses the
already proved rank-one $E$-module structure of rational homology.
The complex structure commutes with $E$ and hence, on the free
rank-one $E\otimes\mathbf R$-module, belongs to $E\otimes\mathbf R$.
Its square is $-1$, which is impossible in a real factor of that
algebra. Therefore $E$ has no real embeddings.

The Rosati involution is nontrivial: for an identity involution the
positive trace form $a\mapsto\operatorname{Tr}_{E/\mathbf Q}(a^2)$
would force every archimedean factor to be real. Indeed a complex
factor has a vector with negative square trace, and rational
coordinates are dense in the real scalar extension, contradicting
positivity. For the fixed field $E^+$ of the nontrivial involution,
the same argument applied to the restricted trace form shows that
$E^+$ is totally real. Since $[E:E^+]=2$ and $E$ is totally imaginary,
$E$ is a CM field, as claimed.

\paragraph{Actual curve and Prym maps.}
For a smooth proper curve $\overline Y/L$, choose a point after a finite
extension and use its Abel--Jacobi morphism
\[
 a:\overline Y\longrightarrow J_{\overline Y},\qquad
 y\longmapsto[y-y_0].
\]
The composite $b=q\circ a:\overline Y\to B$ is nonconstant whenever
$q$ is nonzero: the image of $a$ generates the Jacobian, so a constant
$q\circ a$ (necessarily zero at $y_0$) would imply $q=0$.
The induced homology map is exactly $V_\ell(q)$ under the Abel--Jacobi
isomorphism, and therefore detects the specified line.

In the A.22 finite-cover setting, let
$f:\overline Y\to E_0$ have degree $n$ and let
$P=\ker(f_*:J_{\overline Y}\to E_0)^0$.
The explicit homomorphism
\[
 v=n\,\mathrm{id}_{J_{\overline Y}}-f^*f_*
\]
has image in $P$, since $f_*f^*=[n]$ and its image is connected.
On a line $W\subset\ker V_\ell(f_*)$, it acts by multiplication by
$n$, so it is nonzero on $W$. Apply the preceding argument to the
line $V_\ell(v)(W)\subset V_\ell(P)$. It gives an actual quotient
$q_P:P\to B$ of CM type and an actual nonconstant morphism
\[
 \overline Y\xrightarrow{a}J_{\overline Y}
 \xrightarrow{v}P\xrightarrow{q_P}B
\]
which detects $W$. This construction proves the precise CM-factor
consequence for the residual weight-one branch; it does not merely
attach a representation-theoretic name to the line.

\paragraph{The remaining morphism obligations.}
For $Y=\overline Y\setminus T$, the displayed map $b$ restricts to
$Y\to B\setminus\{0\}$ exactly when $b^{-1}(0)\subset T$, by the
definition of inverse image of the open complement. That containment
has not been established for an arbitrary finite cover in A.22.
Moreover the rank calculation allows CM abelian varieties of dimension
greater than one; it does not produce one of the punctured elliptic
targets covered by F.3.15. F.3.16 supplies both the elliptic targets and
the necessary fibre containment for its particular Fermat cover by
explicit equations. An application to all A.22 covers still requires
the corresponding actual maps and a theorem applying to those targets.

\paragraph{Source status.}
The complete original Faltings1983 PDF was reopened through the primary
Drinfeld archive on19November2026; \S5,S\"atze3--4,p.360 and
Korollar1,p.361 were read directly. The agent's earlier shorthand
message attributing semisimplicity to Korollar1 and Hom to Korollar2
was corrected: Korollar2,p.361 is the isogeny-equivalence statement.
The rank argument above is supplied independently here.
